% =============================================================================
% PREFACIO — Qbit Lab: Laboratorio de Física Computacional
% =============================================================================
\chapter*{Prefacio}
\addcontentsline{toc}{chapter}{Prefacio}

\begin{flushright}
\itshape
``The purpose of computing is insight, not numbers.''\\[2pt]
\upshape --- Richard W. Hamming (1962)
\end{flushright}
\bigskip

Este laboratorio existe porque Hamming tenía razón. El propósito de cada
práctica no es que el estudiante ejecute código y obtenga números, sino que
construya intuición sobre la estructura que las fluctuaciones revelan, los
modos que la regularización selecciona, y las distribuciones que la dinámica
estocástica encuentra. Los números son verificación; la física es el objetivo.

Las prácticas implementan los resultados de \textit{La máquina estocástica}
--- un texto de aprendizaje automático para físicos de posgrado que
sostiene que el ML moderno es física estadística en dimensión alta. Pero
este compendio es un documento autónomo: cada práctica es autocontenida,
con sus propios fundamentos teóricos, y puede seguirse sin haber leído
el libro. Lo que sí se hereda es la perspectiva: aquí no se ``entrena un
modelo''; se deja que un sistema estocástico se relaje hacia el equilibrio,
se observan las consecuencias del principio de máxima entropía, y se
verifica que las ecuaciones de Fokker-Planck y Langevin gobiernan tanto
la difusión de neutrones como el descenso de gradiente estocástico.

\subsection*{Qué se espera del estudiante}

Estas prácticas están dirigidas a estudiantes graduados de física con
formación estándar en mecánica clásica, mecánica estadística, métodos
matemáticos y programación científica. No se asume experiencia previa en
aprendizaje automático. Se asume que el estudiante sabe programar en algún
lenguaje científico y que puede leer una ecuación diferencial sin
sobresaltarse.

Las implementaciones son agnósticas del lenguaje: los algoritmos se
presentan como pseudocódigo y el estudiante los traduce al entorno de su
preferencia (Python, Julia, C++, MATLAB). Se busca comprensión
algorítmica, no dependencia de librerías. Cuando un framework como
PyTorch o JAX aparece, es para verificar resultados, no para reemplazar
el entendimiento.

\subsection*{Estructura de las prácticas}

Cada práctica sigue el formato IMRaD (Introducción, Métodos, Resultados,
Discusión) con secciones para fundamentos teóricos, experimentos numéricos,
análisis de resultados y preguntas abiertas. El arco temático sigue el
del libro:

\begin{enumerate}
  \item Aproximación de funciones con bases fijas: monomios, Fourier, RBF.
        El oscilador armónico como primer laboratorio.
  \item Sobreajuste y regularización: SVD, filtro de Wiener, curva L.
        El condicionamiento como observable numérico.
  \item Geometría del paisaje de pérdida: gradiente, hessiano, saddle points.
        Descenso de gradiente como integrador de Euler.
  \item La red neuronal como aproximador: forward pass, activaciones,
        profundidad vs.\ anchura.
  \item Optimización estocástica: SGD como Langevin, distribución de Gibbs,
        temperatura efectiva y recocido.
  \item Retropropagación: backward pass manual, verificación por diferencias
        finitas, diferenciación automática.
  \item Inferencia y modelos generativos: MLE, MAP, VAE, MC Dropout,
        cuantificación de incertidumbre.
  \item Arquitecturas: convoluciones, atención, transformers.
  \item Ecuaciones diferenciales estocásticas: Euler-Maruyama, Milstein,
        proceso de Ornstein-Uhlenbeck.
  \item Modelos de difusión y flow matching: DDPM, DDIM, score matching,
        transporte óptimo.
\end{enumerate}

Un ejemplo unificado --- la señal de un oscilador armónico simple ---
recorre las primeras siete prácticas como hilo conductor, permitiendo
al estudiante comparar cada herramienta nueva contra las anteriores en
un sistema que conoce de memoria.

\subsection*{Sobre este proyecto}

Este compendio nace como extensión práctica de mi libro \textit{La máquina
estocástica}, un proyecto personal desarrollado durante mis estudios de
doctorado en Ciencias de la Computación. Es también un puente hacia
\texttt{@QbitLab}, mi comunidad de divulgación científica a la que no he
podido retribuir en los últimos años pero que espero revitalizar pronto.
A los más de 40\,000 seguidores que hicieron posible aquella etapa: este
laboratorio es también para ustedes, con la esperanza de que el espíritu
de compartir conocimiento siga vivo en estas páginas.

\vspace{0.3cm}
\noindent
Construido sobre la plantilla \texttt{latex-catppuccin-labs}
(\url{https://github.com/Zessinthel/latex-catppuccin-labs}).

\vspace{1.2cm}
\begin{flushright}
\textit{J.\,A. Chala Casanova\enspace(\texttt{@QbitLab})}\\[2pt]
\the\year
\end{flushright}